\documentclass[11pt,a4paper]{article}
\usepackage[margin=1in]{geometry}
\usepackage{kotex}
\usepackage{amsmath,amssymb,amsfonts,bm,mathtools}
\usepackage{booktabs,multirow,array}
\usepackage{enumitem}
\usepackage{graphicx}
\usepackage{hyperref}
\usepackage{natbib}
\usepackage{xcolor}
\usepackage{setspace}
\usepackage{caption}
\usepackage{float}
\hypersetup{colorlinks=true, linkcolor=blue, citecolor=blue, urlcolor=blue}
\setstretch{1.15}

\title{베이지안 지역 통계 서명과 FiLM-조건부 GRU decoder-only 모형을 결합한\\
태양광 발전량 예측 프레임워크: 모델 설계, 수식화, 기대효과 및 고려사항}
\author{}
\date{}

\begin{document}
\maketitle

\section{연구 목적과 핵심 아이디어}
본 연구의 목적은 지역별 기후, 지형, 구름 동역학, 설비 특성이 상이한 다지역 태양광(PV) 예측 문제에서, \textbf{(i) 베이지안 계층모형이 먼저 지역별 구조적 효과를 추정하고}, \textbf{(ii) FiLM-조건부 GRU decoder-only 모형이 남은 비선형 잔차와 단기 동학을 학습}하도록 하는 2단계 예측 프레임워크를 구축하는 데 있다. 대표적 적용 시나리오는 강원도, 제주도, 경기도처럼 지역별 기후대와 일사 조건이 명확히 다른 경우이다.

핵심 아이디어는 다음과 같다.
\begin{enumerate}[label=(\arabic*)]
    \item \textbf{1단계 베이지안 계층모형}에서 지역별 평균 수준, 기상 민감도, 시간대/연중 계절구조를 posterior 분포로 추정한다.
    \item \textbf{2단계 딥러닝 모형}에서는 1단계가 설명하지 못한 잔차를 대상으로, GRU decoder-only 모형이 비선형성, 시차 의존성, 급변동 패턴을 학습한다.
    \item \textbf{FiLM(feature-wise linear modulation)}은 지역 메타데이터와 1단계 posterior에서 추출한 \emph{regional statistical signature}를 이용하여 공유된 딥러닝 backbone의 내부 표현을 지역별로 조절한다.
\end{enumerate}

최종 예측값은 다음과 같이 정의한다.
\begin{equation}
    \hat y_{i,t+h} = \hat \mu^{\mathrm{Bayes}}_{i,t+h} + \hat e^{\mathrm{DL}}_{i,t+h}, \qquad h=1,\dots,H,
\end{equation}
여기서 $\hat \mu^{\mathrm{Bayes}}_{i,t+h}$는 1단계 베이지안 계층모형의 posterior predictive mean, $\hat e^{\mathrm{DL}}_{i,t+h}$는 2단계 FiLM-조건부 GRU decoder-only 모형이 예측한 잔차이다.

\section{문제 설정과 표기법}
상위 그룹은 `site_cluster`로 고정한다. `site_cluster` 집합을 $\mathcal R = \{1,\dots,R\}$라 하고, 사이트-클러스터 매핑 함수를 $c(i)\in\mathcal R$로 둔다. 클러스터 $r\in\mathcal R$에 속한 사이트 집합을 $\mathcal S_r=\{i:c(i)=r\}$라 하며, 사이트 $i\in\mathcal S_r$의 시점 $t$에서의 관측 발전량을 $y_{i,t}$라 한다. 사이트 $i$의 설치용량을 $C_i$라 두고, 정규화된 출력
\begin{equation}
    \tilde y_{i,t} = \frac{y_{i,t}}{C_i}
\end{equation}
를 사용한다. 이상적으로 $\tilde y_{i,t}\in[0,1]$이다.

입력 변수는 다음과 같이 구분한다.
\begin{itemize}[leftmargin=2em]
    \item \textbf{과거 관측 시계열} $\bm y^{\mathrm{hist}}_{i,t-L+1:t}$: 최근 $L$개 시점의 발전량
    \item \textbf{과거/미래 기상 변수} $\bm x_{i,\tau}$: GHI, 기온, 상대습도, 풍속 (DNI/DHI/운량은 가용 시 파생 또는 보조)
    \item \textbf{달력/태양기하 변수} $\bm c_{\tau}$: hour, day-of-year, weekday, holiday, solar zenith/azimuth
    \item \textbf{클러스터 메타데이터} $\bm m_r$: 클러스터 중심 위경도/고도 통계, 해안성 비율, 변동성 지표, 대표 설비 특성 등
    \item \textbf{품질 변수} $q_{i,\tau}$: `quality_mask`, `dqf_class`, `source_flag` (학습 샘플 필터/가중치 용도, 예측 feature 직접 투입 금지)
\end{itemize}

예측 목적은 시점 $t$에서 미래 $H$개 시점의 벡터
\begin{equation}
    \hat{\bm y}_{i,t+1:t+H} = (\hat y_{i,t+1},\dots,\hat y_{i,t+H})
\end{equation}
를 산출하는 것이다.

\section{1단계: 지역 구조를 추정하는 베이지안 계층모형}
\subsection{설계 원칙}
1단계는 단순 선형회귀가 아니라, \textbf{지역별 구조적 차이와 시간/계절 구조를 posterior 분포 형태로 추정하는 베이지안 계층모형}으로 설계한다. 이 단계의 목적은 단지 baseline을 만드는 것이 아니라, \textbf{지역별 기상 민감도와 시간 구조를 해석 가능한 형태로 추정하고, 그 결과를 2단계의 조건 벡터로 넘기는 것}이다. 베이지안 계층모형은 부분풀링(partial pooling)을 통해 데이터가 적은 지역의 추정을 안정화하고, 동시에 posterior uncertainty까지 정량화할 수 있다는 장점이 있다 \citep{StanGuide,SaintDrenan2019Bayesian}.

\subsection{주간/야간 분리와 관측모형}
PV 출력은 야간에 구조적으로 0이므로, 태양고도 또는 solar zenith angle을 이용한 주간 지시변수
\begin{equation}
    d_t = \mathbb I(\mathrm{SZA}_t < 90^\circ)
\end{equation}
를 정의한다. 그러면
\begin{equation}
    \tilde y_{i,t}=0 \qquad \text{if } d_t=0
\end{equation}
로 두고, 주간 구간($d_t=1$)에서만 확률모형을 적합한다.

주간 정규화 출력이 $(0,1)$ 범위에 있으므로, 베타 회귀(beta regression)를 다음과 같이 사용한다.
\begin{equation}
    \tilde y^{\star}_{i,t} \mid \mu_{i,t},\phi_r \sim \mathrm{Beta}(\mu_{i,t}\phi_r, (1-\mu_{i,t})\phi_r),
    \qquad 0<\mu_{i,t}<1,
    \label{eq:beta_obs}
\end{equation}
여기서 $\tilde y^{\star}_{i,t}$는 $0$과 $1$을 피하기 위해 작은 $\epsilon>0$ (예: $\epsilon=10^{-4}$)에 대해
\begin{equation}
    \tilde y^{\star}_{i,t}=\min\{\max(\tilde y_{i,t},\epsilon),1-\epsilon\}
\end{equation}
로 변환한 값이다. $\phi_r$는 그룹별 정밀도(precision) 파라미터이다.

\subsection{선형예측자: 지역별 계수 + 시간/계절 기저함수}
주간 구간에서 평균 파라미터 $\mu_{i,t}$에 대해 logit link를 사용한다.
\begin{equation}
    \mathrm{logit}(\mu_{i,t}) = \eta_{i,t}.
\end{equation}

본 연구는 GAMM의 smooth 대신, \textbf{베이지안 계층회귀 + Fourier basis expansion}으로 시간대/계절 구조를 직접 모델링한다. 사이트 $i$가 지역 $r(i)$에 속한다고 할 때,
\begin{align}
    \eta_{i,t}
    &= \alpha_0 + a_{r(i)} + u_i
    + \bigl(\beta_{\mathrm{GHI}} + b_{r(i),\mathrm{GHI}}\bigr)\,\widetilde{\mathrm{GHI}}_{i,t}
    + \bigl(\beta_T + b_{r(i),T}\bigr)\,\widetilde T_{i,t}
    + \bigl(\beta_{\mathrm{CLD}} + b_{r(i),\mathrm{CLD}}\bigr)\,\widetilde{\mathrm{CLD}}_{i,t} \\
    &\quad + \beta_{\mathrm{RH}}\widetilde{\mathrm{RH}}_{i,t}
    + \beta_{\mathrm{WS}}\widetilde{\mathrm{WS}}_{i,t}
    + \beta_{\mathrm{CSI}}\widetilde{\mathrm{CSI}}_{i,t}
    + \bm \beta_m^\top \widetilde{\bm m}_i
    + \sum_{k=1}^{K_h}
    \Bigl[\gamma_k^{(s)} \sin\Bigl(\frac{2\pi k h_t}{24}\Bigr)
    + \gamma_k^{(c)} \cos\Bigl(\frac{2\pi k h_t}{24}\Bigr)\Bigr] \\
    &\quad + \sum_{m=1}^{K_d}
    \Bigl[\delta_m^{(s)} \sin\Bigl(\frac{2\pi m d_t^{\mathrm{yr}}}{365}\Bigr)
    + \delta_m^{(c)} \cos\Bigl(\frac{2\pi m d_t^{\mathrm{yr}}}{365}\Bigr)\Bigr] \\
    &\quad + \sum_{k=1}^{K_h}
    \Bigl[\tilde\gamma_{r(i),k}^{(s)} \sin\Bigl(\frac{2\pi k h_t}{24}\Bigr)
    + \tilde\gamma_{r(i),k}^{(c)} \cos\Bigl(\frac{2\pi k h_t}{24}\Bigr)\Bigr] \\
    &\quad + \sum_{m=1}^{K_d}
    \Bigl[\tilde\delta_{r(i),m}^{(s)} \sin\Bigl(\frac{2\pi m d_t^{\mathrm{yr}}}{365}\Bigr)
    + \tilde\delta_{r(i),m}^{(c)} \cos\Bigl(\frac{2\pi m d_t^{\mathrm{yr}}}{365}\Bigr)\Bigr].
    \label{eq:bayes_stage1}
\end{align}

여기서
\begin{itemize}[leftmargin=2em]
    \item $a_{r(i)}$: 지역별 random intercept,
    \item $u_i$: 사이트별 random intercept,
    \item $b_{r(i),\mathrm{GHI}}, b_{r(i),T}, b_{r(i),\mathrm{CLD}}$: 지역별 random slope,
    \item $h_t$: 시간대(hour),
    \item $d_t^{\mathrm{yr}}$: 연중 일자(day of year),
    \item $K_h, K_d$: 각각 일중/연중 Fourier basis 차수,
    \item $\widetilde{\bm m}_i$: site metadata의 표준화 벡터이다.
\end{itemize}

즉, 식 \eqref{eq:bayes_stage1}은
\begin{quote}
\textit{전체 공통 효과 + 지역별 편차 + 사이트별 편차 + 전역 시간/계절 구조 + 지역별 시간/계절 편차}
\end{quote}
를 하나의 베이지안 계층구조로 묶는다.

\subsection{사전분포(priors)}
전역 회귀계수에 대해서는 약정보 사전분포(weakly informative prior)를 둔다.
\begin{equation}
    \alpha_0,\; \beta_{\mathrm{GHI}},\; \beta_T,\; \beta_{\mathrm{CLD}},\; \beta_{\mathrm{RH}},\; \beta_{\mathrm{WS}},\; \beta_{\mathrm{CSI}} \sim \mathcal N(0, 2.5^2).
\end{equation}

지역별 랜덤계수 벡터를
\begin{equation}
    \bm b_r =
    \begin{bmatrix}
        a_r & b_{r,\mathrm{GHI}} & b_{r,T} & b_{r,\mathrm{CLD}}
    \end{bmatrix}^{\top}
\end{equation}
로 두면,
\begin{equation}
    \bm b_r \sim \mathcal N(\bm 0, \bm \Sigma_b),
\end{equation}
이며 공분산행렬은
\begin{equation}
    \bm \Sigma_b = \bm D_b \bm \Omega_b \bm D_b,
\end{equation}
\begin{equation}
    \bm D_b = \mathrm{diag}(\sigma_a,\sigma_{\mathrm{GHI}},\sigma_T,\sigma_{\mathrm{CLD}}),
    \qquad \bm \Omega_b \sim \mathrm{LKJ}(\eta_{\Omega}).
\end{equation}

사이트별 절편은
\begin{equation}
    u_i \sim \mathcal N(0,\sigma_u^2), \qquad \sigma_u \sim \mathrm{half\mbox{-}t}(3,0,1).
\end{equation}

지역별 Fourier 편차계수는 차수별 shrinkage prior를 둔다.
\begin{align}
    \tilde\gamma_{r,k}^{(s)},\tilde\gamma_{r,k}^{(c)} &\sim \mathcal N(0, \tau_{h,k}^2), \\
    \tilde\delta_{r,m}^{(s)},\tilde\delta_{r,m}^{(c)} &\sim \mathcal N(0, \tau_{d,m}^2),
\end{align}
\begin{equation}
    \tau_{h,k},\tau_{d,m} \sim \mathrm{half\mbox{-}t}(3,0,0.5).
\end{equation}
이는 고차 Fourier 성분이 과도하게 커지는 것을 억제한다.

지역별 precision은
\begin{equation}
    \log \phi_r \sim \mathcal N(\mu_\phi,\sigma_\phi^2)
\end{equation}
로 두어 그룹마다 주간 출력 분산 수준이 다를 수 있게 한다. 그룹 수가 매우 작을 때는 $\phi_r$ 대신 사이트 계층 정밀도 $\phi_i$에 상위 하이퍼프라이어를 두는 대안을 우선 고려한다.

\subsection{Posterior predictive baseline과 잔차 정의}
학습 데이터 $\mathcal D_{\mathrm{train}}$에 대해 posterior는
\begin{equation}
    p(\Theta \mid \mathcal D_{\mathrm{train}})
\end{equation}
로 주어진다. 미래 시점의 baseline 예측은 posterior predictive mean으로 정의한다.
\begin{equation}
    \hat \mu^{\mathrm{Bayes}}_{i,t+h}
    = \mathbb E\bigl[\tilde y_{i,t+h}\mid \mathcal D_{\mathrm{train}}\bigr]
    = \int \mathbb E\bigl[\tilde y_{i,t+h}\mid \Theta\bigr] p(\Theta\mid\mathcal D_{\mathrm{train}}) d\Theta.
\end{equation}
MCMC 표본 $\{\Theta^{(s)}\}_{s=1}^S$를 사용하면,
\begin{equation}
    \hat \mu^{\mathrm{Bayes}}_{i,t+h} \approx \frac{1}{S} \sum_{s=1}^S \mu_{i,t+h}^{(s)}.
    \label{eq:posterior_mean}
\end{equation}

이때 1단계 잔차는 원래 발전량 스케일 또는 정규화 스케일에서
\begin{equation}
    e_{i,t} = \tilde y_{i,t} - \hat \mu^{\mathrm{Bayes}}_{i,t}
\end{equation}
로 정의한다. 필요하면 다시 $C_i$를 곱해 원스케일로 복원할 수 있다.

\subsection{Regional statistical signature의 베이지안 정의}
본 연구의 핵심은 지역별 posterior를 그대로 FiLM의 조건벡터로 넘기지 않고, \textbf{posterior에서 추출한 요약 통계량을 지역 서명으로 구성}하는 데 있다. 지역 $r$에 대한 regional statistical signature를 다음과 같이 정의한다.
\begin{equation}
    \bm z_r^{\mathrm{stat}} =
    \Big[
    \underbrace{\mathbb E[a_r\mid\mathcal D],\; \mathbb E[b_{r,\mathrm{GHI}}\mid\mathcal D],\; \mathbb E[b_{r,T}\mid\mathcal D],\; \mathbb E[b_{r,\mathrm{CLD}}\mid\mathcal D]}_{\text{지역 평균 및 민감도}},
    \underbrace{\mathbb E[\tilde\gamma_{r,1}^{(s)}\mid\mathcal D],\dots,\mathbb E[\tilde\delta_{r,K_d}^{(c)}\mid\mathcal D]}_{\text{시간/계절 구조}},
    \underbrace{\mathrm{sd}(a_r\mid\mathcal D),\mathrm{sd}(b_{r,\mathrm{GHI}}\mid\mathcal D),\dots}_{\text{불확실성 요약}}
    \Big].
    \label{eq:regional_signature}
\end{equation}

실무적으로는 식 \eqref{eq:regional_signature}의 차원이 커질 수 있으므로,
\begin{equation}
    \bar{\bm z}_r^{\mathrm{stat}} = \bm P \bm z_r^{\mathrm{stat}}
\end{equation}
와 같은 선형 투영, PCA, 또는 작은 MLP encoder로 압축한 뒤 FiLM generator에 넣는 것이 바람직하다.

\paragraph{선택적 확장: posterior uncertainty의 직접 활용.}
본 연구는 $\bm z_r^{\mathrm{stat}}$에 posterior mean뿐 아니라 posterior standard deviation을 포함할 수 있다. 이는 데이터가 적은 지역일수록 신뢰도가 낮다는 정보를 FiLM이 간접적으로 활용하게 하므로, \textbf{uncertainty-aware conditioning}이라는 추가 novelty로 확장될 수 있다.


\section{2단계: FiLM-조건부 GRU decoder-only 모형에 의한 잔차 예측}
\subsection{왜 잔차를 예측하는가}
1단계가 설명하는 것은 지역별 평균 수준, 기상 민감도, 시간대/연중 계절성 같은 \textbf{구조적 패턴}이다. 그러나 실제 PV 출력은 구름 이동, NWP 오차, 지역별 미세기상, 지형-기상 상호작용, ramp event 등으로 인해 비선형 잔차가 크게 남는다. 따라서 2단계는 원시 발전량을 다시 처음부터 예측하는 대신,
\begin{equation}
    e_{i,t} = \tilde y_{i,t} - \hat \mu^{\mathrm{Bayes}}_{i,t}
\end{equation}
를 학습 대상으로 두는 \textbf{residual forecasting} 구조가 더 타당하다.

\subsection{왜 decoder-only 구조를 사용하는가}
본 연구에서 2단계의 목적은 전체 시계열을 새로 인코딩하여 생성하는 것이 아니라, \textbf{1단계 베이지안 baseline이 설명하지 못한 잔차만 보정하는 것}이다. 따라서 과거 길이 $L$의 residual/history window와 미래 horizon $H$에 대해 이용 가능한 외생변수만으로 충분한 경우, 별도의 encoder--decoder보다 \textbf{decoder-only 예측기}가 더 간결하고 해석 가능하다. 이 구조에서는 하나의 GRU backbone이 과거 residual sequence를 순차적으로 요약하고, 마지막 hidden representation으로부터 $H$-step residual vector를 직접 출력한다. 즉, 예측 대상은
\begin{equation}
    \hat{\bm e}_{i,t+1:t+H} =
    \big(
    \hat e_{i,t+1},\dots,\hat e_{i,t+H}
    \big)
\end{equation}
와 같은 다시계점 residual 벡터이다.

\subsection{딥러닝 입력 정의}
2단계의 과거 입력 시퀀스를 다음과 같이 정의한다.
\begin{equation}
    \bm p_{i,\tau} = \Big[
    e_{i,\tau},\;
    \bm x^{\mathrm{past}}_{i,\tau},\;
    \bm c_{\tau},\;
    \hat \mu^{\mathrm{Bayes}}_{i,\tau}
    \Big],
    \qquad \tau=t-L+1,\dots,t.
\end{equation}
또한 미래 horizon에서 사용 가능한 외생변수는
\begin{equation}
    \bm q_{i,t+h} = \Big[
    \bm x^{\mathrm{fut}}_{i,t+h},\;
    \bm c_{t+h},\;
    \hat \mu^{\mathrm{Bayes}}_{i,t+h}
    \Big],
    \qquad h=1,\dots,H
\end{equation}
로 정의한다. 이때 $\bm x^{\mathrm{fut}}_{i,t+h}$는 반드시 시점 $t$에서 이용 가능한 예보 스냅샷(as-of forecast)으로 구성해 train-inference 불일치를 막는다. 실무적으로는 direct multi-horizon 예측을 위해
\begin{equation}
    \bm q^{\mathrm{stack}}_{i,t} =
    \big[
    \bm q_{i,t+1};\bm q_{i,t+2};\cdots;\bm q_{i,t+H}
    \big]
\end{equation}
처럼 horizon별 미래 정보를 하나의 벡터로 연결하여 출력 head에 공급할 수 있다.

\subsection{Decoder-only GRU backbone}
decoder-only GRU는 과거 입력 시퀀스 $\bm p_{i,t-L+1},\dots,\bm p_{i,t}$를 순차적으로 읽어 최종 hidden state를 형성한다.
\begin{align}
    \bm z_{\tau} &= \sigma(\bm W_z \bm p_{i,\tau} + \bm U_z \bm h_{\tau-1} + \bm b_z), \\
    \bm r_{\tau} &= \sigma(\bm W_r \bm p_{i,\tau} + \bm U_r \bm h_{\tau-1} + \bm b_r), \\
    \tilde{\bm h}_{\tau} &= \tanh\!\Big(\bm W_h \bm p_{i,\tau} + \bm U_h (\bm r_{\tau}\odot \bm h_{\tau-1}) + \bm b_h\Big), \\
    \bm h_{\tau} &= (1-\bm z_{\tau})\odot \bm h_{\tau-1} + \bm z_{\tau}\odot \tilde{\bm h}_{\tau},
\end{align}
\begin{equation}
    \tau=t-L+1,\dots,t.
\end{equation}
최종 hidden state를 $\bm h_{i,t}^{\mathrm{dec}}$라 두면, 이는 과거 residual, 과거 기상, 달력 구조, 베이지안 baseline history를 요약한 표현이다.

\subsection{지역 조건부 FiLM 생성기}
FiLM은 조건 벡터로부터 hidden representation의 feature-wise scale과 shift를 생성한다 \citep{Perez2017FiLM}. 본 연구에서는 조건 벡터를 다음과 같이 둔다.
\begin{equation}
    \bm g_r = \Big[
    \mathrm{Emb}(r),\;
    \bm m_r,\;
    \bar{\bm z}_r^{\mathrm{stat}}
    \Big],
\end{equation}
여기서 $\mathrm{Emb}(r)$는 `site_cluster` ID embedding, $\bm m_r$는 클러스터 메타데이터, $\bar{\bm z}_r^{\mathrm{stat}}$는 압축된 베이지안 클러스터 통계 서명이다.

이 조건 벡터를 작은 MLP에 통과시켜 decoder backbone과 출력 head에서 사용할 FiLM 파라미터를 생성한다.
\begin{equation}
    [\bm \gamma_r^{(\ell)},\;\bm \beta_r^{(\ell)}] = f_{\mathrm{FiLM}}^{(\ell)}(\bm g_r),
\end{equation}
여기서 $\ell$은 modulation이 적용되는 블록을 뜻한다.

\subsection{Decoder-only hidden modulation}
가장 간단한 적용은 최종 hidden state에 대해
\begin{equation}
    \widetilde{\bm h}_{i,t}^{\mathrm{dec}} =
    \bm \gamma_r^{\mathrm{dec}} \odot \bm h_{i,t}^{\mathrm{dec}}
    +
    \bm \beta_r^{\mathrm{dec}}
    \label{eq:film_decoder_only}
\end{equation}
를 적용하는 것이다. 필요하면 중간 GRU hidden에도 같은 방식의 modulation을 둘 수 있지만, 표본 수가 많지 않은 경우에는 식 \eqref{eq:film_decoder_only}와 같은 \textbf{final-state modulation}이 더 안정적이다. 이 식은 동일 backbone의 내부 요약 표현을 지역별로 재보정하는 역할을 한다.

\subsection{Direct multi-horizon residual head}
수정된 hidden state와 미래 외생변수 스택을 결합하여 direct multi-horizon residual vector를 출력한다.
\begin{equation}
    \bm u_{i,t} = \Big[
    \widetilde{\bm h}_{i,t}^{\mathrm{dec}};\;
    \bm q^{\mathrm{stack}}_{i,t}
    \Big]
\end{equation}
그리고 작은 MLP 또는 선형 head를 통해
\begin{equation}
    \hat{\bm e}_{i,t+1:t+H}
    =
    f_{\mathrm{out}}(\bm u_{i,t})
\end{equation}
를 계산한다. 성분별로 쓰면
\begin{equation}
    \hat e_{i,t+h}
    =
    f_{\mathrm{out}}^{(h)}(\bm u_{i,t}),
    \qquad h=1,\dots,H.
\end{equation}
즉, 본 구조는 auto-regressive encoder--decoder처럼 horizon을 순차 생성하지 않고, \textbf{하나의 decoder-only backbone으로부터 미래 $H$개 잔차를 직접 산출하는 구조}이다.

최종 예측값은
\begin{equation}
    \hat{\tilde y}_{i,t+h} = d_{t+h}\cdot \min\{\max(\hat \mu^{\mathrm{Bayes}}_{i,t+h} + \hat e_{i,t+h},0),1\}
    \label{eq:final_forecast}
\end{equation}
로 계산하고,
\begin{equation}
    \hat y_{i,t+h} = C_i \hat{\tilde y}_{i,t+h}
\end{equation}
로 원스케일에 복원한다. 필요하면
\begin{equation}
    \hat{\tilde y}_{i,t+h}^{\mathrm{clip}} = \min\{\max(\hat{\tilde y}_{i,t+h},0),1\}
\end{equation}
로 clipping을 적용한다.

\section{학습 목적함수와 학습 절차}
\subsection{기본 손실함수}
다시계점 잔차 예측의 기본 손실은 품질가중치를 포함한 weighted Huber loss 또는 weighted MSE를 권장한다.
\begin{equation}
    \mathcal L_{\mathrm{res}} = \sum_{h=1}^H w_h \, q_{i,t+h}\, d_{t+h}\,\ell\Big(e_{i,t+h}, \hat e_{i,t+h}\Big),
\end{equation}
여기서 $w_h$는 horizon별 가중치이며 가까운 horizon에 더 큰 값을 둘 수 있다. $q_{i,t+h}\in[0,1]$는 품질가중치(`quality_mask` 기반)이다.

최종적으로는 다음과 같은 복합 손실을 사용할 수 있다.
\begin{equation}
    \mathcal L = \lambda_1 \mathcal L_{\mathrm{res}} + \lambda_2 \mathcal L_{\mathrm{range}} + \lambda_3 \mathcal L_{\mathrm{agg}} + \lambda_4 \lVert \Theta_{\mathrm{FiLM}} \rVert_2^2.
\end{equation}

각 항의 의미는 다음과 같다.
\begin{align}
    \mathcal L_{\mathrm{range}} &= \sum_{h=1}^H \Big[\max(0,-\hat{\tilde y}_{i,t+h})^2 + \max(0,\hat{\tilde y}_{i,t+h}-1)^2\Big], \\
    \mathcal L_{\mathrm{agg}} &= \sum_{r\in\mathcal R}\sum_{h=1}^H \Big(Y^{\mathrm{obs}}_{r,t+h} - \sum_{i\in\mathcal S_r} \hat y_{i,t+h}\Big)^2.
\end{align}

즉, $\mathcal L_{\mathrm{range}}$는 물리적 범위 제약, $\mathcal L_{\mathrm{agg}}$는 지역 집계 정합성을 위한 항이다.

\subsection{누수 방지를 위한 2단계 학습 절차}
본 연구에서는 정보누수를 방지하기 위해 다음 절차를 권장한다.
\begin{enumerate}[label=Step \arabic*., leftmargin=2.5em]
    \item 시간 순서를 유지한 train/validation/test split 또는 rolling-origin evaluation을 구성한다.
    \item 각 학습 fold에서 \textbf{1단계 베이지안 계층모형을 해당 fold의 training 구간에만 적합}한다.
    \item 같은 training 구간 내부에서도 out-of-fold 또는 nested rolling 방식으로 $\hat \mu^{\mathrm{Bayes}}$와 $e_t$를 생성하여 2단계 입력을 만든다.
    \item 2단계 GRU decoder-only + FiLM 모형은 오직 1단계에서 생성된 baseline/residual 및 이용 가능한 미래 외생변수만 사용하여 학습한다.
    \item 테스트 시점에서는 training data로 적합된 1단계 posterior predictive mean을 생성한 후, 이를 2단계에 입력하여 최종 예측을 수행한다.
\end{enumerate}

이 절차를 지키지 않으면 1단계 baseline이 미래 정보를 암묵적으로 사용하게 되어, 2단계 성능이 과대평가될 수 있다.

\section{기대효과 및 고려점}
\subsection{기대효과}
본 프레임워크의 기대효과는 다음과 같다.
\begin{enumerate}[label=(\arabic*)]
    \item \textbf{구조와 비구조의 명시적 분리}: 1단계가 지역별 구조적 패턴을 설명하고, 2단계가 잔차 동학에 집중함으로써 학습 부담을 분산시킨다.
    \item \textbf{저표본 지역 일반화}: 베이지안 부분풀링과 공유 backbone 덕분에, 데이터가 적은 지역에서도 완전 독립모형보다 안정적인 추정이 가능하다.
    \item \textbf{해석 가능성}: 단순 region one-hot과 달리, 베이지안 통계 서명 $\bm z_r^{\mathrm{stat}}$는 ``이 지역이 어떤 기상 민감도와 시간/계절 구조를 가지는가''를 posterior 기반으로 제공한다.
    \item \textbf{불확실성 정량화}: 1단계 posterior variance를 통해 지역별 추정 신뢰도를 수치화할 수 있으며, 필요하면 이를 FiLM 조건벡터에 직접 포함할 수 있다.
    \item \textbf{공유 표현 + 지역 적응}: backbone은 공유하여 표본 효율성을 확보하고, FiLM은 지역 조건부 modulation을 통해 지역 특화 반응을 부여한다.
    \item \textbf{다시계점 예측 적합성}: direct multi-horizon decoder-only 구조를 통해 $1$-step부터 $H$-step까지 horizon-dependent error pattern을 학습할 수 있다.
\end{enumerate}

\subsection{고려점과 잠재적 한계}
반면 다음 사항은 반드시 고려해야 한다.
\begin{enumerate}[label=(\alph*)]
    \item \textbf{지역 수가 너무 적은 경우}: 강원/제주/경기처럼 지역 카테고리만 3개이고 사이트 수도 적다면 FiLM 자유도가 과해질 수 있다. 이 경우 각 지역 내 다수 사이트를 확보하거나, FiLM generator를 작게 설계해야 한다.
    \item \textbf{베이지안 1단계의 계산비용}: MCMC 기반 계층모형은 계산비용이 크다. 실무에서는 NUTS 기반 학습 후 posterior sample 수를 제한하거나, variational inference를 보조적으로 사용할 수 있다.
    \item \textbf{야간 0 처리}: 베타 회귀는 $(0,1)$ 구간에서 정의되므로, 야간 deterministic zero 분리 또는 zero-one-inflated beta 구조를 명확히 정해야 한다.
    \item \textbf{1단계가 너무 약한 경우}: baseline이 부정확하면 residual이 구조적 패턴을 다시 포함하게 되어, 2단계가 사실상 원시 문제를 다시 풀게 된다.
    \item \textbf{1단계가 너무 강한 경우}: 반대로 baseline이 거의 모든 분산을 설명하면 2단계가 학습할 정보가 부족해질 수 있다. 그러나 이는 해석성 측면에서는 오히려 장점일 수도 있다.
    \item \textbf{데이터 누수}: 2단계 입력으로 사용하는 posterior predictive baseline 및 residual 생성 시 반드시 out-of-fold 원칙을 지켜야 한다.
    \item \textbf{지역 메타데이터 품질}: 위경도, 고도, 설비 orientation 등 메타정보 품질이 낮으면 FiLM conditioning 효과가 제한될 수 있다.
\end{enumerate}

\section{Novelty를 어디에 둘 것인가}
본 연구의 novelty는 단순히 ``GRU에 지역 정보를 넣었다''가 아니라, 다음 다섯 가지 결합점에 두는 것이 바람직하다.

\subsection{Novelty 1: 지역 ID가 아니라 posterior 기반 \emph{통계적 지역 서명}을 조건으로 사용}
기존 다지역 PV 예측은 지역 ID embedding, 위치 좌표, 또는 site embedding을 활용하는 경우가 많다 \citep{Goutte2024Embedding,Wen2024DecoderGeo}. 그러나 본 연구는 1단계 베이지안 계층모형이 추정한 $\bm z_r^{\mathrm{stat}}$를 FiLM 조건벡터에 포함함으로써, 단순 범주형 식별자를 넘어선 \textbf{해석 가능하고 posterior로 요약된 지역 표현}을 사용한다는 점이 차별적이다.

\subsection{Novelty 2: 구조적 지역효과와 비선형 잔차동학의 명시적 분해}
regional PV forecasting 연구는 통계적 upscaling/parameterization 계열과 딥러닝 계열로 분리되어 발전해 왔다 \citep{SaintDrenan2019Bayesian,Jung2022VAR,Perera2024HTCNN}. 본 연구는 이 둘을 직렬로 연결하여, \textbf{베이지안 통계모형이 구조를 식별하고 딥러닝이 구조 밖의 잔차를 학습}하도록 설계한다. 이 점이 method-level novelty의 핵심이다.

\subsection{Novelty 3: posterior uncertainty까지 활용 가능한 conditioning}
기존 location conditioning은 주로 point representation에 머무는 경우가 많다. 본 연구는 posterior mean뿐 아니라 posterior standard deviation 또는 credible interval 기반 요약치를 $\bm z_r^{\mathrm{stat}}$에 포함할 수 있으므로, \textbf{uncertainty-aware regional conditioning}으로 확장할 수 있다. 이는 특히 저표본 지역이나 기상 변동성이 큰 지역에서 의미가 크다.

\subsection{Novelty 4: 공유 backbone + 지역별 modulation의 결합}
기존 multi-site/common-model 접근은 공통 모델이 여러 사이트를 동시에 학습하도록 설계되지만, 지역별 동질성과 이질성을 동시에 다루는 방식은 제한적이다 \citep{Jang2024CommonModel}. 본 연구는 backbone 파라미터는 공유하되, FiLM을 통해 \textbf{같은 특징 채널이라도 지역별로 다른 scale/shift를 부여}한다. 따라서 표본 효율성과 지역 적응성을 동시에 추구할 수 있다.

\subsection{Novelty 5: 저표본 지역과 일반화 평가까지 포함한 방법론 프레임}
최근 리뷰들은 PV forecasting에서 benchmark와 평가 통일성이 부족함을 지적한다 \citep{Iheanetu2022Review}. 따라서 본 연구는 모델 구조뿐 아니라, \textbf{지역별/계절별/horizon별 성능, 저표본 지역 성능, unseen site 일반화}까지 포함한 평가 프레임워크를 제안하는 방향으로 novelty를 확장할 수 있다.

\section{논문 본문에 바로 쓸 수 있는 기여 문장 예시}
본 연구의 기여는 다음과 같이 정리할 수 있다.
\begin{enumerate}[label=\textbf{C\arabic*.}, leftmargin=2.2em]
    \item 지역별 구조적 효과를 베이지안 계층모형으로 추정하고, 이로부터 posterior 기반의 해석 가능한 regional statistical signature를 구성하는 2단계 PV 예측 프레임워크를 제안한다.
    \item regional statistical signature와 지역 메타데이터를 FiLM conditioning에 결합함으로써, 공유된 GRU decoder-only backbone의 내부 표현을 지역별로 조절하는 방법을 제안한다.
    \item 구조적 baseline과 비선형 residual forecasting을 명시적으로 분해하여, 다지역 PV 예측에서 정확도와 해석 가능성을 동시에 확보하는 방법을 제시한다.
    \item posterior uncertainty를 지역 서명에 포함할 수 있는 확장 구조를 제안하여, 저표본 지역 또는 기상 변동성이 큰 지역에서 uncertainty-aware conditioning 가능성을 제시한다.
    \item 지역별, 계절별, horizon별, 저표본 조건에서의 ablation 및 일반화 실험을 통해 제안 방법의 효과를 검증한다.
\end{enumerate}

\section{실험 설계 권장안}
논문화 가능성을 높이려면 최소한 다음 비교군이 필요하다.
\begin{table}[H]
\centering
\caption{권장 비교군(ablation) 구성}
\begin{tabular}{ll}
\toprule
구분 & 설명 \\
\midrule
B1 & Bayesian hierarchical model only \\
B2 & GRU only \\
B3 & GRU + region one-hot / embedding \\
B4 & GRU + FiLM (without Bayesian signature) \\
B5 & Bayesian baseline + GRU residual \\
B6 & Bayesian baseline + GRU + FiLM (metadata only) \\
B7 & Bayesian baseline + GRU + FiLM + posterior regional signature (제안모형) \\
B8 & B7 + posterior uncertainty features (선택 확장) \\
\bottomrule
\end{tabular}
\end{table}

평가 지표는 MAE, RMSE, nMAE, nRMSE를 기본으로 하고, 지역별/계절별/horizon별 성능을 반드시 분리 보고하는 것이 좋다. 또한 제주처럼 기상 변동이 큰 지역, 강원처럼 지형성이 강한 지역, 경기처럼 발전 패턴이 다른 지역에서 제안모형의 상대 향상을 강조하는 것이 바람직하다.

\section{요약}
정리하면, 본 연구는 \textbf{``베이지안 지역 통계 서명에 의해 조건부 조절되는 잔차 기반 딥러닝 PV 예측''}이라는 프레임으로 정리할 수 있다. 1단계 베이지안 계층모형은 지역별 구조를 posterior 형태로 해석 가능하게 식별하고, 2단계 FiLM-조건부 GRU decoder-only 모형은 지역별 비선형 잔차 동학을 학습한다. 따라서 본 방법은 단순 지역 ID 주입 방식보다 더 설명력 있고, 완전 독립 지역모형보다 더 표본 효율적이며, 순수 블랙박스 딥러닝보다 더 해석 가능한 지역 적응형 PV 예측 프레임워크가 될 수 있다.

\begin{thebibliography}{99}

\bibitem[Birnbaum et~al.(2019)Birnbaum, Kuleshov, Enam, Koh, and Ermon]{Birnbaum2019TFiLM}
Birnbaum, S., Kuleshov, V., Enam, Z., Koh, P. W., and Ermon, S. (2019).
\newblock Temporal FiLM: Capturing Long-Range Sequence Dependencies with Feature-Wise Modulations.
\newblock In \emph{Advances in Neural Information Processing Systems}.

\bibitem[Goutte et~al.(2024)Goutte, Klotzner, and Le]{Goutte2024Embedding}
Goutte, S., Klotzner, K., and Le, H.-V. (2024).
\newblock Forecasting photovoltaic production with neural networks and weather features.
\newblock \emph{Energy Economics}, 139, 107884.

\bibitem[Iheanetu et~al.(2022)Iheanetu, Idemudia, and Aliu]{Iheanetu2022Review}
Iheanetu, K. J., Idemudia, C. U., and Aliu, S. O. (2022).
\newblock Solar Photovoltaic Power Forecasting: A Review.
\newblock \emph{Sustainability}, 14(24), 17005.

\bibitem[Jang et~al.(2024)Jang, Lee, and coauthors]{Jang2024CommonModel}
Jang, S. Y., Lee, D., et al. (2024).
\newblock A Deep Learning-Based Solar Power Generation Forecasting Model for Multiple Sites Using a Single Common Model.
\newblock \emph{Sustainability}, 16(12), 5240.

\bibitem[Jung et~al.(2022)Jung, Lee, Kim, Kim, Kim, and Lee]{Jung2022VAR}
Jung, A. H., Lee, D. H., Kim, J. Y., Kim, C. K., Kim, H. G., and Lee, Y. S. (2022).
\newblock Regional Photovoltaic Power Forecasting Using Vector Autoregression Model in South Korea.
\newblock \emph{Energies}, 15(21), 7853.

\bibitem[Perez et~al.(2017)Perez, Strub, de~Vries, Dumoulin, and Courville]{Perez2017FiLM}
Perez, E., Strub, F., de~Vries, H., Dumoulin, V., and Courville, A. (2017).
\newblock FiLM: Visual Reasoning with a General Conditioning Layer.
\newblock arXiv preprint arXiv:1709.07871.

\bibitem[Perera et~al.(2024)Perera, and coauthors]{Perera2024HTCNN}
Perera, M., et al. (2024).
\newblock Day-ahead regional solar power forecasting with hierarchical temporal convolutional neural networks using historical power generation and weather data.
\newblock \emph{Applied Energy}, 370, 123310.

\bibitem[Saint-Drenan et~al.(2019)Saint-Drenan, Vogt, Killinger, Bright, and Potthast]{SaintDrenan2019Bayesian}
Saint-Drenan, Y.-M., Vogt, S., Killinger, S., Bright, J. M., and Potthast, R. (2019).
\newblock Bayesian parameterisation of a regional photovoltaic model -- Application to forecasting.
\newblock \emph{Solar Energy}, 188, 592--605.

\bibitem[Stan Development Team(2025)]{StanGuide}
Stan Development Team. (2025).
\newblock \emph{Stan User's Guide}.
\newblock \url{https://mc-stan.org/docs/stan-users-guide/}.

\bibitem[Wen et~al.(2024)Wen and coauthors]{Wen2024DecoderGeo}
Wen, Y., et al. (2024).
\newblock Improving multi-site photovoltaic forecasting with relevance amplification.
\newblock \emph{Energy}, 299, 131591.

\bibitem[Yang et~al.(2017)Yang, Quan, Disfani, and Li]{Yang2017GeoHierarchy}
Yang, D., Quan, H., Disfani, V. R., and Li, S. (2017).
\newblock Reconciling solar forecasts: Geographical hierarchy.
\newblock \emph{Solar Energy}, 146, 276--286.

\bibitem[Yagli et~al.(2019)Yagli, Yang, and Srinivasan]{Yagli2019Sequential}
Yagli, G. M., Yang, D., and Srinivasan, D. (2019).
\newblock Reconciling solar forecasts: Sequential reconciliation.
\newblock \emph{Solar Energy}, 191, 391--400.

\end{thebibliography}

\end{document}
